\subsection{Using the Log of a Power to Solve Equations} 
If we know that both sides of an equation are \makeblank[1in], we can take the log of both sides. This works especially well if one or both sides contains an \makeblank[1in] expression.
\begin{example}
Solve $3^x=5$.
\begin{itemize}
\item We cannot easily write 3 and 5 as \makeblank[1in] of the same number.
\item Therefore, we cannot solve this by \makeblank.
\item We'll take the \makeblank of both sides, so we can use a calculator.
\end{itemize}
\begin{lpic}{}{3}
\end{lpic}
\end{example}
We could also solve this by writing the equation in log form:
\[
\makeblank[3in]
\]
This suggests the last property of logarithms.